\documentclass{beamer}

% --- Thème et options ---
\usetheme{Madrid}
\usecolortheme{default}
\setbeamertemplate{navigation symbols}{}
\setbeamertemplate{footline}[frame number]

% --- Packages ---
\usepackage[T1]{fontenc}
\usepackage[utf8]{inputenc}
\usepackage[french]{babel}
\usepackage{amsmath}
\usepackage{amsfonts}
\usepackage{amssymb}
\usepackage{tikz}
\usetikzlibrary{decorations.pathreplacing, positioning, arrows.meta}
\usepackage{float}
\usepackage{caption}

% --- Définition de l'environnement importantnote pour beamer ---
\newenvironment{importantnote}
{%
    \begin{block}{Note importante}
}
{%
    \end{block}
}

% --- Informations du document ---
\title{Équilibrage des lignes de production}
\subtitle{Line Balancing - Principes fondamentaux}
\author{}
\institute{}
\date{}

\begin{document}

% --- Slide 1: Titre ---
\begin{frame}
    \titlepage
    \begin{center}
        \small Méthodologie Lean Manufacturing
    \end{center}
\end{frame}

% --- Slide 2: Sommaire ---
\begin{frame}{Sommaire}
    \tableofcontents
\end{frame}

% ============================================================
% PARTIE 1: INTRODUCTION À L'ÉQUILIBRAGE
% ============================================================
\section{Introduction à l'équilibrage}

\begin{frame}{Définition}
    \begin{block}{Qu'est-ce que l'équilibrage des lignes ?}
        L'équilibrage des lignes de production (\textit{Line Balancing}) est une méthode fondamentale du Lean Manufacturing qui vise à répartir équitablement la charge de travail entre les différents postes d'une chaîne de production.
    \end{block}
    
    \vspace{0.5cm}
    
    \textbf{Objectifs :}
    \begin{itemize}
        \item Minimiser les temps morts
        \item Optimiser l'utilisation des ressources
        \item Améliorer le flux de production
    \end{itemize}
\end{frame}

\begin{frame}{Objectifs de l'équilibrage}
    \begin{columns}
        \column{0.5\textwidth}
        \begin{itemize}
            \item Chaque poste a une durée proche du \textbf{Takt Time}
            \item Respect des contraintes de précédence
            \item Nombre de postes minimisé
            \item Réduction des temps morts
        \end{itemize}
        
        \column{0.5\textwidth}
        \begin{importantnote}
            Un équilibrage parfait est atteint lorsque la durée de chaque poste est exactement égale au Takt Time, éliminant ainsi toute attente.
        \end{importantnote}
    \end{columns}
\end{frame}

\begin{frame}{Le concept de goulot d'étranglement}
    \begin{definition}[Goulot d'étranglement (Bottleneck)]
        Un poste de travail qui reçoit plus de travail que ce qu'il peut traiter à sa capacité de production maximale.
    \end{definition}
    
    \vspace{0.5cm}
    
    \textbf{Conséquences :}
    \begin{itemize}
        \item Interruption du flux de travail
        \item Retards dans le processus de production
        \item Accumulation de stocks intermédiaires
    \end{itemize}
    
    \vspace{0.3cm}
    \textcolor{blue}{→ L'équilibrage vise à identifier et résoudre ces goulots.}
\end{frame}

% ============================================================
% PARTIE 2: LES CONCEPTS FONDAMENTAUX
% ============================================================
\section{Concepts fondamentaux}

\begin{frame}{Trois concepts clés à distinguer}
    \begin{center}
        \begin{tikzpicture}
            \node[draw, rounded corners, fill=blue!20, minimum width=2.5cm, minimum height=1cm] (takt) at (0,0) {Takt Time};
            \node[draw, rounded corners, fill=green!20, minimum width=2.5cm, minimum height=1cm] (cycle) at (4,0) {Cycle Time};
            \node[draw, rounded corners, fill=orange!20, minimum width=2.5cm, minimum height=1cm] (lead) at (8,0) {Lead Time};
            
            \node[below=0.3cm of takt] {Rythme client};
            \node[below=0.3cm of cycle] {Temps réel};
            \node[below=0.3cm of lead] {Temps total};
        \end{tikzpicture}
    \end{center}
    
    \vspace{0.5cm}
    \begin{itemize}
        \item Souvent confondus mais complémentaires
        \item Chacun joue un rôle spécifique dans l'équilibrage
    \end{itemize}
\end{frame}

\begin{frame}{Takt Time ($T_T$)}
    \begin{block}{Définition}
        Le Takt Time est le rythme de production dicté par la demande client. Il représente le temps maximum autorisé pour produire une unité.
    \end{block}
    
    \begin{center}
        \boxed{T_T = \frac{\text{Temps de production disponible par période}}{\text{Demande client par période}}}
    \end{center}
    
    \vspace{0.3cm}
    \textbf{Exemple :}
    \[
    T_T = \frac{7\text{h} \times 60}{600 \text{ pièces}} = 0,7 \text{ min/pièce} = 42 \text{ secondes/pièce}
    \]
\end{frame}

\begin{frame}{Caractéristiques du Takt Time}
    \begin{columns}
        \column{0.5\textwidth}
        \begin{itemize}
            \item Donnée objective basée sur la demande client
            \item Sert de référence pour l'équilibrage
            \item Varie selon les fluctuations de la demande
            \item Ne dépend pas de la capacité de production
        \end{itemize}
        
        \column{0.5\textwidth}
        \begin{center}
            \begin{tikzpicture}
                \draw[->] (0,0) -- (4,0) node[right] {Temps};
                \draw[dashed, red, thick] (0,1) -- (4,1) node[right] {$T_T$};
                \fill[blue!30] (0.5,0) rectangle (1.5,1);
                \fill[blue!30] (1.8,0) rectangle (2.8,1);
                \fill[blue!30] (3.1,0) rectangle (3.8,1);
                \node at (1,0.5) {Pièce};
                \node at (2.3,0.5) {Pièce};
                \node at (3.45,0.5) {Pièce};
            \end{tikzpicture}
            \small{Rythme imposé par le client}
        \end{center}
    \end{columns}
\end{frame}

\begin{frame}{Cycle Time ($T_C$)}
    \begin{block}{Définition}
        Le Cycle Time est le temps réel entre deux pièces produites sur un poste donné.
    \end{block}
    
    \begin{center}
        \boxed{T_C = \frac{\text{Temps de production d'une pièce sur un poste}}{\text{Nombre de pièces produites}}}
    \end{center}
    
    \vspace{0.3cm}
    \textbf{Interprétation :}
    \begin{itemize}
        \item $T_C < T_T$ : Surcapacité (temps morts)
        \item $T_C > T_T$ : Goulot (incapacité à suivre la demande)
        \item $T_C = T_T$ : Équilibre parfait
    \end{itemize}
\end{frame}

\begin{frame}{Lead Time ($LT$)}
    \begin{block}{Définition}
        Le Lead Time est le temps total entre l'entrée de la matière première et la sortie du produit fini.
    \end{block}
    
    \begin{center}
        \boxed{LT = \sum (\text{Temps de transformation}) + \sum (\text{Temps d'attente}) + \sum (\text{Temps de transport})}
    \end{center}
    
    \vspace{0.3cm}
    \textbf{Relation importante :}
    \[
    LT = N \times T_C + \text{temps d'attente}
    \]
    où $N$ est le nombre de postes.
\end{frame}

\begin{frame}{Tableau récapitulatif}
    \begin{table}[H]
    \centering
    \begin{tabular}{|p{2.2cm}|p{3cm}|p{3cm}|p{3cm}|}
    \hline
    \textbf{Concept} & \textbf{Définition} & \textbf{Formule} & \textbf{Rôle} \\
    \hline
    Takt Time & Rythme imposé par le client & $T_T = \frac{T_{\text{dispo}}}{D}$ & Objectif à atteindre \\
    \hline
    Cycle Time & Temps réel de production & $T_C = \frac{T_{\text{prod}}}{Q}$ & Indicateur d'écart \\
    \hline
    Lead Time & Temps total de traversée & $LT = \sum T_C + \sum \text{attente}$ & Indicateur global \\
    \hline
    \end{tabular}
    \end{table}
\end{frame}

\begin{frame}{Illustration graphique}
    \begin{center}
    \begin{tikzpicture}[scale=0.7, every node/.style={scale=0.7}]
        % Timeline
        \draw[thick, ->] (0,0) -- (12,0) node[right] {Temps};
        
        % Lead Time
        \draw[decorate, decoration={brace, amplitude=10pt, mirror}] (0.5,-0.8) -- (11.5,-0.8) node[midway, below] {Lead Time};
        
        % Postes - CORRECTION : utilisation du point pour les décimales
        \fill[blue!20] (1,0) rectangle (3,0.6) node[midway, above] {Poste 1\\$T_C=0.62$};
        \fill[blue!20] (3.5,0) rectangle (5.5,0.6) node[midway, above] {Poste 2\\$T_C=0.66$};
        \fill[blue!20] (6,0) rectangle (8,0.6) node[midway, above] {Poste 3\\$T_C=0.70$};
        \fill[blue!20] (8.5,0) rectangle (10.5,0.6) node[midway, above] {Poste 4\\$T_C=0.63$};
        
        % Attentes
        \draw[<->, red] (3,0.3) -- (3.5,0.3) node[midway, above] {Attente};
        \draw[<->, red] (5.5,0.3) -- (6,0.3) node[midway, above] {Attente};
        \draw[<->, red] (8,0.3) -- (8.5,0.3) node[midway, above] {Attente};
        
        % Takt Time line
        \draw[dashed, green!60!black, thick] (0,1) -- (12,1) node[right] {Takt Time = 0.7 min};
    \end{tikzpicture}
    \end{center}
    \vspace{0.3cm}
    \begin{center}
        \textcolor{blue}{Le Cycle Time du poste le plus lent détermine le rythme de production}
    \end{center}
\end{frame}

% ============================================================
% PARTIE 3: POURQUOI ÉQUILIBRER ?
% ============================================================
\section{Pourquoi équilibrer ?}

\begin{frame}{Les bénéfices d'un bon équilibrage}
    \begin{columns}
        \column{0.5\textwidth}
        \begin{itemize}
            \item \textbf{Réduction des temps morts}
            \item \textbf{Amélioration du flux de production}
            \item \textbf{Optimisation des ressources}
        \end{itemize}
        
        \column{0.5\textwidth}
        \begin{itemize}
            \item \textbf{Réduction des stocks WIP}
            \item \textbf{Augmentation de la productivité}
            \item \textbf{Meilleure visibilité des problèmes}
        \end{itemize}
    \end{columns}
    
    \vspace{0.5cm}
    \begin{center}
        \textcolor{green}{\textbf{→ Gain de compétitivité globale}}
    \end{center}
\end{frame}

\begin{frame}{Les coûts du déséquilibre}
    \begin{table}[H]
    \centering
    \begin{tabular}{|p{3cm}|p{7cm}|}
    \hline
    \textbf{Type de perte} & \textbf{Conséquences} \\
    \hline
    Temps morts & Sous-utilisation des opérateurs et des machines \\
    \hline
    Stocks WIP & Capital immobilisé, risque d'obsolescence, surfaces occupées \\
    \hline
    Lead time long & Manque de réactivité, délais clients allongés \\
    \hline
    Goulots & Production limitée par le poste le plus lent \\
    \hline
    Surcapacité & Investissements sous-optimisés \\
    \hline
    \end{tabular}
    \end{table}
\end{frame}

\begin{frame}{Impact sur la performance globale}
    \begin{center}
    \begin{tikzpicture}[node distance=1cm, scale=0.9, every node/.style={scale=0.9}]
        \node[draw, rectangle, fill=blue!10, minimum width=3cm, minimum height=0.8cm] (lean) at (0,0) {Lean Manufacturing};
        \node[draw, rectangle, fill=green!10, minimum width=3cm, minimum height=0.8cm, below of=lean] (equi) {Équilibrage};
        \node[draw, rectangle, fill=yellow!10, minimum width=2.5cm, minimum height=0.8cm, below left of=equi, xshift=-1.5cm] (prod) {Productivité};
        \node[draw, rectangle, fill=yellow!10, minimum width=2.5cm, minimum height=0.8cm, below of=equi] (qual) {Qualité};
        \node[draw, rectangle, fill=yellow!10, minimum width=2.5cm, minimum height=0.8cm, below right of=equi, xshift=1.5cm] (delai) {Délais};
        \node[draw, rectangle, fill=red!10, minimum width=4cm, minimum height=0.8cm, below of=prod, xshift=2cm] (client) {Satisfaction Client};
        
        \draw[->, thick] (lean) -- (equi);
        \draw[->, thick] (equi) -- (prod);
        \draw[->, thick] (equi) -- (qual);
        \draw[->, thick] (equi) -- (delai);
        \draw[->, thick] (prod) -- (client);
        \draw[->, thick] (qual) -- (client);
        \draw[->, thick] (delai) -- (client);
    \end{tikzpicture}
    \end{center}
\end{frame}

% ============================================================
% PARTIE 4: QUAND ET COMMENT ?
% ============================================================
\section{Quand et comment équilibrer ?}

\begin{frame}{Quand faut-il équilibrer ?}
    \begin{columns}
        \column{0.5\textwidth}
        \begin{itemize}
            \item Création d'une nouvelle ligne
            \item Changement significatif de la demande
            \item Introduction d'un nouveau produit
        \end{itemize}
        
        \column{0.5\textwidth}
        \begin{itemize}
            \item Identification de goulots récurrents
            \item Taux d'efficacité < 80\%
            \item Stocks intermédiaires anormalement élevés
        \end{itemize}
    \end{columns}
    
    \vspace{0.5cm}
    \begin{importantnote}
        L'équilibrage n'est pas une action ponctuelle mais une démarche d'amélioration continue. Les lignes doivent être régulièrement rééquilibrées.
    \end{importantnote}
\end{frame}

\begin{frame}{Résumé}
    \begin{center}
        \begin{tikzpicture}
            \node[draw, rounded corners, fill=blue!20, minimum width=8cm, minimum height=0.8cm] (titre) at (0,0) {\textbf{L'équilibrage des lignes de production}};
            
            \node[draw, rounded corners, fill=green!10, minimum width=3.5cm, minimum height=2.5cm, below left=0.5cm of titre] (def) {\begin{tabular}{c} \textbf{Définition} \\ Répartition \\ équitable des \\ charges \end{tabular}};
            
            \node[draw, rounded corners, fill=yellow!10, minimum width=3.5cm, minimum height=2.5cm, below=0.5cm of titre] (obj) {\begin{tabular}{c} \textbf{Objectif} \\ Atteindre \\ le Takt Time \end{tabular}};
            
            \node[draw, rounded corners, fill=orange!10, minimum width=3.5cm, minimum height=2.5cm, below right=0.5cm of titre] (ben) {\begin{tabular}{c} \textbf{Bénéfices} \\ Productivité \\ Qualité \\ Délais \end{tabular}};
        \end{tikzpicture}
    \end{center}
\end{frame}

\begin{frame}{Conclusion}
    \begin{center}
        \Large
        \textcolor{blue}{« L'équilibrage est le cœur du Lean Manufacturing »}
    \end{center}
    
    \vspace{1cm}
    
    \begin{itemize}
        \item Un outil fondamental pour l'optimisation industrielle
        \item Une approche structurée basée sur des concepts clés
        \item Des bénéfices mesurables sur la performance globale
        \item Une démarche d'amélioration continue
    \end{itemize}
    
    \vspace{1cm}
    \begin{center}
        \textbf{À retenir :} Le Takt Time est le rythme, le Cycle Time est la réalité, l'équilibrage est la solution.
    \end{center}
\end{frame}

\end{document}
